\documentclass[11pt]{article}
\usepackage[T1]{fontenc}
\usepackage{lmodern}
\usepackage[margin=1in,headheight=15pt]{geometry}
\usepackage{amsmath,amssymb,amsthm,mathtools}
\usepackage{booktabs,array,longtable,microtype,xurl,fancyhdr}
\usepackage[dvipsnames]{xcolor}
\usepackage[colorlinks=true,linkcolor=MidnightBlue,citecolor=MidnightBlue,urlcolor=MidnightBlue]{hyperref}
\usepackage{bookmark,needspace}
\numberwithin{equation}{section}
\newtheorem{theorem}{Theorem}[section]
\newtheorem{lemma}[theorem]{Lemma}
\newtheorem{proposition}[theorem]{Proposition}
\newtheorem{corollary}[theorem]{Corollary}
\newtheorem{conjecture}[theorem]{Conjecture}
\theoremstyle{definition}\newtheorem{definition}[theorem]{Definition}
\theoremstyle{remark}\newtheorem{remark}[theorem]{Remark}
\DeclareMathOperator{\Li}{Li}
\DeclareMathOperator{\Cov}{Cov}
\DeclareMathOperator{\supp}{supp}
\newcommand{\dd}{\,\mathrm d}
\newcommand{\ii}{\mathrm i}
\newcommand{\E}{\mathbb E}
\newcommand{\R}{\mathbb R}
\newcommand{\C}{\mathbb C}
\newcommand{\Q}{\mathbb Q}
\newcommand{\N}{\mathbb N}
\newcommand{\cut}{\C\setminus[1,\infty)}
\newcommand{\HH}{\mathcal H}
\newcommand{\BB}{\mathrm B}
\newcommand{\CC}{C_*}
\newcommand{\pin}{\nolinkurl{5537940fde730713eaac752265a486832b2fe540}}
\emergencystretch=3em
\setlength{\parskip}{3pt}
\pagestyle{fancy}\fancyhf{}
\fancyhead[L]{\small Beta transfer and critical Euler errors}
\fancyhead[R]{\small ProveIt}
\fancyfoot[C]{\thepage}
\hypersetup{pdftitle={Beta Transfer for Harmonic Polylogarithms},pdfauthor={ProveIt Contributors},pdfsubject={Proof of the sharp critical Euler constant, positive densities, uniform asymptotics and cyclotomic identities}}
\title{\textbf{Beta Transfer for Harmonic Polylogarithms}\\[9pt]
\large A sharp critical Euler constant, positive remainder densities,\\
uniform error asymptotics, and cyclotomic profile identities}
\author{ProveIt Contributors}
\date{9 October 2026, America/Los\_Angeles\\[3pt]
\small Research continuation at repository revision \pin}
\begin{document}
\maketitle\thispagestyle{empty}
\begin{abstract}
For $a,b>0$, put $F_{a,b}(z)=\sum_{n\ge1}H_{n-1}^{(b)}z^n/n^a$,
continued to $\C\setminus[1,\infty)$. A beta-distribution representation
shows that, at fixed $w=a+b$, the function
$a\mapsto-\Im F_{a,w-a}(\ii\rho)$ is strictly decreasing for
$0<\rho\le\sqrt3$. Its critical specialization proves Conjecture 12.3 of
\emph{The Sharp Real-Order Threshold for Harmonic Polylogarithms}:
the best constant uniform over $a\in(0,1)$ and Euler truncations $N\ge1$ is
$\pi/4+(\log2)/2$.

A positive divided-difference remainder extends one-sided Euler
evaluation to all positive orders, including $a+b<1$, where the finite
signed Hausdorff measure used in earlier proofs does not exist. On the
critical line a new elementary beta-averaged density satisfies the
parameter-uniform bounds $1<\nu_a(T)<1+1/T$. This yields strengthened
harmonic-defect inequalities, a completely monotone Hurwitz remainder,
and a full positive-coefficient large-$T$ expansion. Uniform averaging
of that expansion gives a moving worst-parameter theorem:
if $\ell_N=\tfrac12\log N$, then the maximal scaled Euler error is
$I_N\{1+(1+o(1))/(e\ell_N\log\ell_N)\}$, and every maximizing parameter satisfies
$(1-a_N)\log\ell_N\to1$.

Rational points of the beta profile admit finite cyclotomic
polylogarithm reductions, with necessary lower-order terms at repeated
poles. Exact rational certificates cover twelve Gaussian values,
including three below the moment threshold, and disprove the stronger
claim that every scaled error decreases with $a$: it already fails at
$N=8$. Six cyclotomic reductions and 176 symbolic equalities are replayed
separately from numerical diagnostics. The proofs are mathematical,
not proof-assistant formalizations; the existing $S_6$ identity and the
broader angular-zero conjectures remain unresolved here.
\end{abstract}
\clearpage
\begingroup\small
\setlength{\parskip}{0pt}
\tableofcontents
\endgroup
\clearpage

\section{The question resolved and the scope of the continuation}
\label{sec:scope}

\subsection{A precise conjecture rather than a new numerical fit}
The repository baseline is \pin. Its canonical polylogarithm manuscript
and incoming-package ledger distinguish integrated proofs from historical
experiments \cite{repo}. In particular, $S_4$ is already proved, and the
existing $S_6$ reduction remains conjectural. Neither status is reset by
this article.

The immediate predecessor is the preserved report \emph{The Sharp
Real-Order Threshold for Harmonic Polylogarithms} \cite{threshold}. Its
Conjecture 12.3, source label \texttt{conj:constant}, states that
\begin{equation}\label{eq:oldconj}
 a\longmapsto-2\Im F_{a,1-a}(\ii)
 \quad\hbox{decreases strictly on }(0,1),
\end{equation}
and proposes the sharp critical Euler constant
\begin{equation}\label{eq:Cstar}
 \CC=\frac\pi4+\frac{\log2}{2}
 =1.13197175367742096\ldots.
\end{equation}
The predecessor proves the endpoint limits and the fixed-parameter
interpretation of this constant, but explicitly leaves its monotonicity
and uniform upper bound open. We prove both. More importantly, the proof
works on every positive total-order line $a+b=w$, not merely at $w=1$.

\subsection{What is prior, and what is added here}
The positive double integral for $F_{a,b}$, the critical endpoint atom,
and the finite signed-measure threshold $a+b\ge1$ are prior results in
\cite{threshold}. We rederive the integral and the critical measure to
make the new arguments independently readable. The beta transfer,
pointwise density bounds, the all-positive-order Euler remainder proof,
the uniform worst-parameter asymptotics, and the cyclotomic profile
formula are the principal additions.

\begin{center}\small
\begin{tabular}{@{}p{.24\textwidth}p{.31\textwidth}p{.36\textwidth}@{}}
\toprule
Topic & Inspected predecessor & Present result\\
\midrule
Critical Euler constant & Conjecture 12.3 & Strict monotonicity and sharp constant\\[3pt]
Order comparison & Critical endpoint data & Transfer theorem at every $w>0$\\[3pt]
Euler evaluation & Signed-measure proof for $w\ge1$ & Positive remainder for all $a,b>0$\\[3pt]
Critical density & Hurwitz plus beta regularizer & Elementary positive excess; uniform bounds\\[3pt]
Large truncations & Sharp $N$-dependence proposed & Full density expansion; uniform maximizing scale\\[3pt]
Finite reductions & General cyclotomic framework & Rational-profile formula including repeated poles\\
\bottomrule
\end{tabular}
\end{center}

A targeted audit of these sources is not an exhaustive priority search.
Gamma and beta integral calculus, stochastic ordering of beta densities,
partial fractions, and Laplace asymptotics are classical tools. The
article claims the proved statements and their consequences relative to
the inspected research program, not that every special case is absent
from all previous literature. General positive Hausdorff-moment theory
is discussed in Liu and Pego \cite{liupego}; the signed and parameter-varying
structures here require their own proofs.

\subsection{Conventions and boundary values}
For $a,b>0$, define
\begin{equation}\label{eq:Fdef}
 H_m^{(b)}=\sum_{j=1}^m j^{-b},\qquad
 c_n(a,b)=\frac{H_{n-1}^{(b)}}{n^a},\qquad
 F_{a,b}(z)=\sum_{n=1}^\infty c_n(a,b)z^n\quad(|z|<1).
\end{equation}
Thus $F_{a,b}(z)=\Li_{a,b}(z,1)$ in the nested-series convention
$\Li_{a,b}(x,y)=\sum_{n>m\ge1}x^ny^m/(n^am^b)$.
Set $g_{a,b}=\Im F_{a,b}(\ii)$, where the value is the continuation
constructed below. All logarithms of positive real variables are real;
complex polylogarithms use the principal slit-plane continuation
\cite{dlmfpolylog}. The entire functions $\eta(s)$ and $\beta(s)$ denote
the Dirichlet eta and beta functions, continued from their defining
alternating series where necessary.

The distinction between an Abel value and an ordinary boundary sum is
essential. If $0<b<1$, integral comparison gives
\begin{equation}\label{eq:coefficientgrowth}
 c_n(a,b)\sim\frac{n^{1-a-b}}{1-b}.
\end{equation}
On $a+b=1$ the coefficients approach $1/a$; below that line they grow.
Consequently the odd-index alternating series defining $g_{a,b}$ does
not converge ordinarily in either case. The Euler evaluator proved
below computes the analytic, equivalently radial Abel, value. It does
not turn the divergent original series into an ordinarily convergent
one.

\section{A beta distribution moves order between the two indices}
\label{sec:beta}

\begin{lemma}[Double gamma integral]\label{lem:double}
For all $a,b>0$ and $z\in\cut$,
\begin{equation}\label{eq:double}
 F_{a,b}(z)=\frac1{\Gamma(a)\Gamma(b)}
 \int_0^\infty\!\int_0^\infty
 s^{a-1}t^{b-1}
 \frac{z^2e^{-2s-t}}
 {(1-ze^{-s})(1-ze^{-s-t})}\dd s\dd t.
\end{equation}
The integral is locally uniformly convergent on the slit plane.
\end{lemma}
\begin{proof}
Inside the disk, insert the gamma integrals for $n^{-a}$ and $m^{-b}$
in the nested sum. Write $n=m+k$, with $m,k\ge1$. The two geometric sums
then give the displayed kernel. Absolute convergence permits every
interchange. For a compact subset of the slit plane,
$|1-zu|$ is bounded below uniformly for $0\le u\le1$. The absolute
kernel is consequently at most a compact-set constant times
$e^{-2s-t}$. This is integrable against the two gamma powers.
Local uniform convergence and holomorphy prove the continuation.
\end{proof}

Put
\begin{equation}\label{eq:kernel}
 K_z(x,y)=\frac{z^2xy}{(1-zx)(1-zy)},\qquad
 \HH_w(v;z)=\frac1{\Gamma(w)}\int_0^\infty
 T^{w-1}K_z(e^{-vT},e^{-T})\dd T,
\end{equation}
for $w>0$ and $0\le v\le1$. We call $\HH_w$ the \emph{beta profile};
this is a definition, not another notation for $F$ at $a=wv$.

\begin{theorem}[Beta transfer identity]\label{thm:beta}
Let $a,b>0$, $w=a+b$, and $V\sim\mathrm{Beta}(a,b)$. Then
\begin{equation}\label{eq:beta}
 F_{a,b}(z)=\E\,\HH_w(V;z)
 =\frac1{\BB(a,b)}\int_0^1v^{a-1}(1-v)^{b-1}\HH_w(v;z)\dd v.
\end{equation}
The profile is continuous on $[0,1]$, locally uniformly in $z\in\cut$,
and its endpoints are
\begin{align}
 \HH_w(0;z)&=\frac z{1-z}\Li_w(z),\label{eq:end0}\\
 \HH_w(1;z)&=\Li_{w-1}(z)-\Li_w(z).\label{eq:end1}
\end{align}
\end{theorem}
\begin{proof}
Use $T=s+t$, $v=s/(s+t)$ in \eqref{eq:double}. The Jacobian is $T$,
and the powers become $T^{w-1}v^{a-1}(1-v)^{b-1}$. Since
$\BB(a,b)=\Gamma(a)\Gamma(b)/\Gamma(w)$, this proves \eqref{eq:beta}.
Uniformly in $v\in[0,1]$, the kernel in \eqref{eq:kernel} is bounded
near $T=0$ and is $O(e^{-T})$ at infinity on compact slit sets. This
proves endpoint continuity by dominated convergence. At $v=0$ the
kernel is $z/(1-z)$ times the classical polylogarithm kernel. At $v=1$
its series coefficient is $(n-1)/n^w$, giving \eqref{eq:end1} first in
the disk and then by analytic continuation.
\end{proof}

\begin{lemma}[Strict beta order comparison]\label{lem:betaorder}
Fix $w>0$. If $J$ is continuous and strictly decreasing on $[0,1]$, then
$a\mapsto\E J(V_a)$ is strictly decreasing for $0<a<w$, where
$V_a\sim\mathrm{Beta}(a,w-a)$. More precisely,
\begin{equation}\label{eq:covariance}
 \frac{\dd}{\dd a}\E J(V_a)
 =\Cov\left(J(V_a),\log\frac{V_a}{1-V_a}\right)<0.
\end{equation}
\end{lemma}
\begin{proof}
On any compact subinterval of $0<a<w$, the beta density and its
$a$-derivative have a common integrable majorant; the derivative adds
only logarithmic endpoint factors. Differentiating the normalized
integral gives \eqref{eq:covariance}. If $V_a'$ is independent with the
same law, the covariance is half the expectation of
\[
 [J(V_a)-J(V_a')]
 \left[\log\frac{V_a}{1-V_a}-\log\frac{V_a'}{1-V_a'}\right].
\]
This integrand is strictly negative off the diagonal, a set of full
product measure. This proves strictness without an appeal to a
numerical parameter grid.
\end{proof}

\section{The order-transfer theorem and the conjectured constant}
\label{sec:monotonicity}

\begin{theorem}[Strict transfer of Gaussian and real values]\label{thm:transfer}
Fix $w>0$. For every $0<\rho\le\sqrt3$,
\begin{equation}\label{eq:transfer}
 a\longmapsto-\Im F_{a,w-a}(\ii\rho)
 \quad\hbox{is strictly decreasing on }(0,w).
\end{equation}
For every real $z<1$ with $z\ne0$, the positive function
$a\mapsto F_{a,w-a}(z)$ is likewise strictly decreasing.
\end{theorem}
\begin{proof}
For $0<y<x<1$,
\begin{equation}\label{eq:G}
 -\Im K_{\ii\rho}(x,y)
 =G_\rho(x,y)
 =\frac{\rho^3xy(x+y)}{(1+\rho^2x^2)(1+\rho^2y^2)}>0.
\end{equation}
At fixed $y$, the sign of its $x$-derivative is the sign of
\begin{equation}\label{eq:Gderiv}
 \frac{\dd}{\dd x}\frac{x(x+y)}{1+\rho^2x^2}
 =\frac{2x+y-\rho^2x^2y}{(1+\rho^2x^2)^2}.
\end{equation}
For $\rho^2\le3$ the numerator is positive: using $x<1$,
$2x+y-3x^2y>2x-2y\ge0$. Thus $G_\rho(x,y)$ increases strictly with
$x$. In \eqref{eq:kernel}, increasing $v$ decreases $x=e^{-vT}$ at
fixed $y=e^{-T}$. Hence $J(v)=-\Im\HH_w(v;\ii\rho)$ decreases
strictly. Apply Lemma~\ref{lem:betaorder}.

For real $z<1$, all denominators are positive, and
\[
 \partial_x K_z(x,y)=\frac{z^2y}{(1-zy)(1-zx)^2}>0.
\]
The same proof applies. No period-independence assertion is involved.
\end{proof}

\begin{remark}
The radius $\sqrt3$ is a sufficient range for a pointwise kernel
argument. This theorem does not claim it is the largest possible
radius for the integrated order comparison. The radial monotonicity
of angular zeros is a different question and is not proved here.
\end{remark}

\begin{corollary}[Sharp fixed-weight endpoint bounds]\label{cor:endbounds}
For $w>0$ and $0<a<w$,
\begin{equation}\label{eq:endbounds}
 \beta(w)-\beta(w-1)
 <-g_{a,w-a}
 <\frac{\beta(w)+2^{-w}\eta(w)}2.
\end{equation}
Both bounds are sharp endpoint limits as $a\uparrow w$ and $a\downarrow0$,
respectively. In particular,
\begin{equation}\label{eq:criticalrange}
 \frac\pi2-1<-2g_{a,1-a}<\CC\qquad(0<a<1),
\end{equation}
and the middle expression decreases strictly with $a$.
\end{corollary}
\begin{proof}
As $a\downarrow0$, $V_a$ converges weakly to the point mass at zero;
as $a\uparrow w$, it converges to the point mass at one. For example,
$\E V_a=a/w$ proves the first limit in probability, and the second
follows by symmetry. Apply the continuity in Theorem~\ref{thm:beta}.
Since
$\Li_w(\ii)=-2^{-w}\eta(w)+\ii\beta(w)$, equations
\eqref{eq:end0}--\eqref{eq:end1} give \eqref{eq:endbounds}.
At $w=1$, use $\beta(1)=\pi/4$, $\beta(0)=1/2$ and $\eta(1)=\log2$.
\end{proof}

This proves the strict-monotonicity part of the predecessor's
Conjecture 12.3. Section~\ref{sec:criticalerror} proves that the same
number is exactly the optimal constant for every critical Euler
truncation, not just for the Gaussian values themselves.

\section{Euler evaluation without a finite signed measure}
\label{sec:eulerall}

Define, for $k\ge0$ and $N\ge1$,
\begin{equation}\label{eq:eulerdef}
 f(n)=c_{2n+1}(a,b),\qquad
 d_k=\sum_{n=0}^k(-1)^n\binom kn f(n),\qquad
 E_N=\sum_{k=0}^{N-1}\frac{d_k}{2^{k+1}}.
\end{equation}
In particular $f(0)=d_0=E_1=0$. Put
\begin{equation}\label{eq:rN}
 r_N(u)=\frac{(1-u^2)^N}{1+u^2},\qquad
 L_N(x,y)=\frac{xy}{x-y}\{r_N(y)-r_N(x)\}.
\end{equation}
Here $L_N(x,y)$ is a two-variable remainder kernel; later the symbol
$\ell_N=\tfrac12\log N$ will denote the scalar logarithmic scale.

\begin{theorem}[One-sided Euler convergence at all positive orders]
\label{thm:allEuler}
For every $a,b>0$, $d_k<0$ for $k\ge1$, and $E_N$ decreases strictly
to $g_{a,b}$. The exact remainder is
\begin{equation}\label{eq:allerror}
 2^N(E_N-g_{a,b})
 =\frac1{\Gamma(a)\Gamma(b)}
 \int_0^\infty\!\int_0^\infty
 s^{a-1}t^{b-1}L_N(e^{-s},e^{-s-t})\dd s\dd t>0.
\end{equation}
For $N\ge2$ let
\[
 m_N=\frac{(2N-2)^{N-1}}{(2N-1)^{N-1/2}},\qquad m_1=1.
\]
Then
\begin{equation}\label{eq:allbound}
 0<E_N-g_{a,b}
 \le 2^{-a}\,2(N+1)m_N\,2^{-N}
 <2(N+1)2^{-N}.
\end{equation}
In particular the sharper displayed bound is $O_a(\sqrt N\,2^{-N})$,
uniform in $b>0$.
\end{theorem}
\begin{proof}
For fixed $x>y>0$, coefficient comparison in
\begin{equation}\label{eq:twoatom}
 K_z(x,y)=\frac{zxy}{x-y}
 \left(\frac1{1-zx}-\frac1{1-zy}\right)
\end{equation}
gives the odd coefficient $xy(x^{2n}-y^{2n})/(x-y)$.
Its $k$th binomial difference is
\[
 \frac{xy}{x-y}\{(1-x^2)^k-(1-y^2)^k\}<0\quad(k\ge1).
\]
Integrate with the positive gamma weight from \eqref{eq:double}.
All sums used here are finite, so there is no unproved summation
interchange.

The finite geometric sum for each of the two terms in
\eqref{eq:twoatom} gives the exact identity
$2^N(E_N-g)=L_N(x,y)$ for the elementary kernel. Integrating proves
\eqref{eq:allerror}. To justify the resulting integral directly, note
that $r_N$ is strictly decreasing and the mean-value theorem gives
\[
 0<L_N(x,y)\le xy\,\|r_N'\|_\infty.
\]
Moreover,
\begin{equation}\label{eq:rderivative}
 |r_N'(u)|
 =2u(1-u^2)^{N-1}
 \left\{\frac N{1+u^2}+\frac{1-u^2}{(1+u^2)^2}\right\}
 \le2(N+1)u(1-u^2)^{N-1}.
\end{equation}
The maximum of the final factor is $m_N$, by a single differentiation
(or $m_1=1$). Finally,
\[
 \frac1{\Gamma(a)\Gamma(b)}\iint
 s^{a-1}t^{b-1}e^{-2s-t}\dd s\dd t=2^{-a}.
\]
This proves the bound and convergence. Strict decrease of $E_N$ follows
from $E_{N+1}-E_N=d_N/2^{N+1}<0$.
\end{proof}

\begin{remark}[Why this extends the earlier domain]
For $a+b<1$, \eqref{eq:coefficientgrowth} precludes a finite signed
measure $\mu$ on $[0,1]$ with $c_n=\int u^{n-1}\dd\mu$: those moments
would be bounded by $\|\mu\|_{\mathrm{TV}}$. Equation
\eqref{eq:allerror} does not assert the existence of such a measure.
It integrates a divided difference whose cancellation is kept intact.
The factor $(x-y)^{-1}$ is controlled by $r_N(y)-r_N(x)$ before
integration. Separating those two terms without justification would
lose the key estimate.
\end{remark}

\begin{proposition}[Finite weights and coefficient stability]\label{prop:weights}
The known finite Euler rearrangement takes the form
\begin{equation}\label{eq:finiteweights}
 E_N=2^{-N}\sum_{n=0}^{N-1}(-1)^nf(n)
                   \sum_{j=n+1}^N\binom Nj.
\end{equation}
The sum of the absolute weights in this expression is $N/2$.
Consequently perturbing every $f(n)$ by at most $\varepsilon$ changes
$E_N$ by at most $N\varepsilon/2$.
\end{proposition}
\begin{proof}
The coefficient identity
$\sum_{k=n}^{N-1}2^{N-k-1}\binom kn=\sum_{j=n+1}^N\binom Nj$
follows from Pascal's rule and gives \eqref{eq:finiteweights}.
Summing its nonnegative absolute weights reverses the order of a
finite sum and yields
$2^{-N}\sum_{j=1}^Nj\binom Nj=N/2$.
\end{proof}
The finite-weight formula itself is already used in the predecessor.
The stability observation explains why its exact-interval implementation
need not suffer the intermediate growth of a full difference table.
The present work does not claim a new bit-complexity theorem for
arbitrary-precision power evaluation.

\section{An elementary positive density on the critical line}
\label{sec:density}
Throughout this section $0<a<1$, $b=1-a$, and $V\sim\mathrm{Beta}(a,b)$.
For $T>0$ and $0<v<1$, define
\begin{equation}\label{eq:Dv}
 D_v(T)=1+\frac1{e^{(1-v)T}-1}
       -\frac1v\frac1{e^{(1-v)T/v}-1},\qquad
 \nu_a(T)=\E D_V(T).
\end{equation}
The notation $\nu_a(T)$ is a density function evaluated at a logarithmic
coordinate. The corresponding measure is
$\nu_a(-\log u)\dd u$, \emph{not} $\nu_a(T)\dd T$.

\begin{lemma}[A parameter-uniform pointwise bound]\label{lem:D}
For every $T>0$ and $0<v<1$,
\begin{equation}\label{eq:Dbound}
 1<D_v(T)<1+\frac1T.
\end{equation}
Moreover, $D_v(T)$ extends continuously to $v=0,1$, with
\begin{equation}\label{eq:Dends}
 D_0(T)=\frac1{1-e^{-T}},\qquad D_1(T)=1.
\end{equation}
\end{lemma}
\begin{proof}
Put $g(x)=x/(e^x-1)$ for $x>0$. It decreases strictly from $1$ to $0$.
For completeness, the sign of $g'$ follows from
$e^x(1-x)-1<0$, whose derivative is $-xe^x$.
It is also strictly convex, because
\[
 g''(x)=\frac{e^x\{(x-2)e^x+x+2\}}{(e^x-1)^3}>0.
\]
Indeed the expression in braces and its first derivative vanish at
zero, while its second derivative is $xe^x>0$.
Thus $h_0=1-g$ is strictly concave, increasing, and $h_0(0)=0$,
$0<h_0<1$.

Set $x=(1-v)T$ and $y=x/v$. Then
\[
 D_v(T)-1=\frac{g(vy)-g(y)}{vy}.
\]
Strict decrease of $g$ proves positivity. Concavity gives
$h_0(vy)>v h_0(y)$, whence
$g(vy)-g(y)<(1-v)h_0(y)<1-v$. Since $vy=(1-v)T$,
this proves the upper bound. The limit at $v=0$ follows directly;
at $v=1$, use $1/(e^x-1)=x^{-1}-1/2+O(x)$ in the difference.
\end{proof}

\begin{theorem}[Critical measure and positive excess]\label{thm:measure}
The critical coefficient sequence has the representation
\begin{equation}\label{eq:measure}
 \mu_a=\frac1a\delta_1-\nu_a(-\log u)\dd u,\qquad
 c_n(a,1-a)=\int_{[0,1]}u^{n-1}\dd\mu_a(u).
\end{equation}
The positive density has mass $1/a$ and satisfies
\begin{equation}\label{eq:nubound}
 1<\nu_a(T)<1+\frac1T\qquad(T>0).
\end{equation}
In particular the atom and negative density in \eqref{eq:measure}
are the Jordan components, and $\|\mu_a\|_{\mathrm{TV}}=2/a$.
\end{theorem}
\begin{proof}
The existence of the atom is known from \cite{threshold}; the following
construction additionally proves the pointwise excess bound.
For $0<v<1$ put
\begin{equation}\label{eq:Psi}
 \Psi_v(u)=\int_{u^{1/v}}^u\frac{\dd y}{1-y^{1-v}}\quad(0<u<1).
\end{equation}
Differentiating its endpoints gives
$\Psi_v'(u)=D_v(-\log u)>0$. Also $\Psi_v(0)=0$ and
\begin{equation}\label{eq:mv}
 \Psi_v(1):=\lim_{u\uparrow1}\Psi_v(u)
 =m_v:=\frac{-\log v}{1-v}.
\end{equation}
To check the last limit, replace the integrand near one by
$[(1-v)(1-y)]^{-1}$; their difference is bounded, while the interval
of integration shrinks. The ratio of its endpoint deficits tends to
$v$, giving \eqref{eq:mv}.

The coefficient of $z^n$ in $\HH_1(v;z)$ is
\begin{equation}\label{eq:profilecoeff}
 h_n(v)=\int_0^1
 \frac{y^{v(n-1)}-y^{n-1}}{1-y^{1-v}}\dd y.
\end{equation}
For $n\ge2$, Tonelli applied to \eqref{eq:Psi} shows that
$(n-1)\int_0^1u^{n-2}\Psi_v(u)\dd u=h_n(v)$: the inner interval
for $u$ is $[y,y^v]$. Integration by parts now gives
\begin{equation}\label{eq:fixedprofilemeasure}
 h_n(v)=m_v-\int_0^1u^{n-1}D_v(-\log u)\dd u.
\end{equation}
For $n=1$ the same equation states $0=m_v-m_v$.

A beta average of $m_v$ can be evaluated exactly:
\begin{equation}\label{eq:betamass}
 \E m_V
 =\frac1{\BB(a,1-a)}\int_0^1
 (-\log v)v^{a-1}(1-v)^{-a-1}\dd v=\frac1a.
\end{equation}
Use
$\frac{\dd}{\dd v}(v/(1-v))^a
 =a v^{a-1}(1-v)^{-a-1}$ and integrate by parts; both endpoint
terms vanish, and the remaining integral is $\BB(a,1-a)/a$.
The positivity in Lemma~\ref{lem:D} and \eqref{eq:betamass} justify
all beta averages by Tonelli. Averaging \eqref{eq:fixedprofilemeasure}
and using Theorem~\ref{thm:beta} proves \eqref{eq:measure}.
The pointwise bound follows by averaging \eqref{eq:Dbound}.
\end{proof}

\begin{corollary}[Strict harmonic-defect moments]\label{cor:defects}
For $n\ge1$, set
\begin{equation}\label{eq:defects}
 r_n(a)=\frac1a-\frac{H_{n-1}^{(1-a)}}{n^a}-\frac1n.
\end{equation}
Then
\begin{equation}\label{eq:defectmoments}
 r_n(a)=\int_0^1u^{n-1}\{\nu_a(-\log u)-1\}\dd u>0.
\end{equation}
Every finite alternating difference is strictly positive:
$\sum_{j=0}^k(-1)^j\binom kjr_{n+j}(a)>0$ for $k\ge0$.
For every $m\ge1$ and integer shift $n\ge1$, the matrix
$(r_{n+i+j}(a))_{0\le i,j<m}$ is positive definite, so its determinant
is strictly positive.
\end{corollary}
\begin{proof}
Subtract the moment of Lebesgue measure from \eqref{eq:measure}.
Finite differences insert the positive factor $(1-u)^k$.
For a nonzero vector $(q_0,\ldots,q_{m-1})$, the associated quadratic
form is the integral of
$u^{n-1}(\sum q_j u^j)^2$ against a density strictly positive on
$(0,1)$. A nonzero polynomial has only finitely many zeros there.
\end{proof}
The strengthened subtraction $1/n$ is justified by
$\nu_a>1$; it is not obtained merely from positivity of the earlier
critical defect.

\subsection{An exact identity with the earlier Hurwitz density}
The new representation can be matched to the regularized Hurwitz
kernel, making the cancellation structure explicit.
Let
\[
 h(t)=\frac1{1-e^{-t}}-\frac1t,\qquad
 C_b(T)=\frac1{\BB(b,a)}\int_0^1 u^{b-1}(1-u)^{a-1}h(Tu)\dd u.
\]
Here $a=1-b$, and the beta parameters of $C_b$ are \emph{reversed}
relative to $\nu_a$.

\begin{proposition}[Beta--Hurwitz density identity]\label{prop:hurwitzdensity}
For $0<a<1$, $b=1-a$ and $T>0$,
\begin{equation}\label{eq:oldnew}
 \E_{V\sim\mathrm{Beta}(a,b)}D_V(T)
 =-\frac{\zeta(b)}{\Gamma(a)}T^{-b}
   +\E_{U\sim\mathrm{Beta}(b,a)}h(TU).
\end{equation}
\end{proposition}
\begin{proof}
The standard regularized Hurwitz integral is
\begin{equation}\label{eq:hurwitzreg}
 \zeta(b,q)=\frac{q^{1-b}}{b-1}
 +\frac1{\Gamma(b)}\int_0^\infty t^{b-1}e^{-qt}h(t)\dd t
 \quad(q>0,\ 0<b<1).
\end{equation}
It follows by subtracting $1/t$ from the classical integral and
continuing in $b$; the subtracted integral is now locally integrable
for $\Re b>0$ \cite{dlmfhurwitz}. The same formula is used in
\cite{threshold}.
Convolving its integral kernel with $T^{a-1}/\Gamma(a)$ and using
$a+b=1$ gives
\[
 \int_0^\infty e^{-qT}C_b(T)\dd T
 =q^{-a}\zeta(b,q)+\frac1a.
\]
Thus the Laplace transform of the right side of \eqref{eq:oldnew}
is $1/a-q^{-a}\{\zeta(b)-\zeta(b,q)\}$. This right-side density
is positive, since $\zeta(b)<0$ and $h>0$, and has mass $1/a$ after
multiplication by $e^{-T}\dd T$. For integer $q=n$, its moments in
$u=e^{-T}$ agree with those of $\nu_a(-\log u)\dd u$ by
\eqref{eq:measure}. Uniqueness follows from density of polynomials in
$C[0,1]$ for finite measures. Both densities are continuous on
$T>0$, so equality almost everywhere implies the displayed pointwise
identity.
\end{proof}

\begin{corollary}[A completely monotone refined Hurwitz remainder]
\label{cor:Hurwitz}
For $q>0$ and $b=1-a\in(0,1)$, put
\begin{equation}\label{eq:Rq}
 R_a(q)=\frac1a-q^{-a}\{\zeta(b)-\zeta(b,q)\}-\frac1q.
\end{equation}
Then
\begin{equation}\label{eq:Rlaplace}
 R_a(q)=\int_0^\infty e^{-qT}\{\nu_a(T)-1\}\dd T>0.
\end{equation}
For all integers $k\ge1$ and $h>0$,
\begin{align}
 0<(-1)^kR_a^{(k)}(q)&<\frac{\Gamma(k)}{q^k},\label{eq:Rderiv}\\
 0<R_a(q)-R_a(q+h)&<\log\left(1+\frac hq\right).\label{eq:Rdiff}
\end{align}
\end{corollary}
\begin{proof}
The Laplace calculation in Proposition~\ref{prop:hurwitzdensity} proves
\eqref{eq:Rlaplace} for every real $q>0$, not merely at integers.
Its convergence near zero follows either from \eqref{eq:oldnew} or
from its positive Laplace construction. Differentiation under the
integral is justified on compact $q$-intervals. For $k\ge1$, the bound
$\nu_a(T)-1<1/T$ gives the gamma integral in \eqref{eq:Rderiv}.
The same bound gives
$\int_0^\infty(e^{-qT}-e^{-(q+h)T})\dd T/T=
\log((q+h)/q)$, proving \eqref{eq:Rdiff}.
\end{proof}

\section{Sharp critical Euler bounds and a stronger claim that fails}
\label{sec:criticalerror}
Define the scaled critical error
\begin{equation}\label{eq:AN}
 A_N(a)=2^N(E_N-g_{a,1-a})\qquad(0<a<1,\ N\ge1),
\end{equation}
and the elementary comparison integrals
\begin{equation}\label{eq:IN}
 I_N=\int_0^1\frac{(1-u^2)^N}{1+u^2}\dd u,\qquad
 B_N=\int_0^1\frac{(1-u^2)^N}{(1+u^2)(1-u)}\dd u\quad(N\ge1).
\end{equation}

\begin{theorem}[The optimal critical constant]\label{thm:sharpconstant}
For every $0<a<1$ and $N\ge1$,
\begin{equation}\label{eq:criticalexact}
 A_N(a)=\int_0^1r_N(u)\nu_a(-\log u)\dd u.
\end{equation}
The sequence $A_N(a)$ decreases strictly with $N$, and
\begin{equation}\label{eq:sharpsandwich}
 I_N<A_N(a)\le A_1(a)=-2g_{a,1-a}<\CC.
\end{equation}
The equality $A_N=A_1$ in the weak inequality occurs only at $N=1$.
Consequently
\begin{equation}\label{eq:sharpEuler}
 0<E_N-g_{a,1-a}<\CC\,2^{-N},
\end{equation}
and $\CC$ is the least constant valid uniformly in $a\in(0,1)$ and
$N\ge1$. The lower comparison $I_N$ is also sharp, as $a\uparrow1$.
\end{theorem}
\begin{proof}
For a finite signed moment measure, summing the finite geometric series
in \eqref{eq:eulerdef} gives
\[
 2^N(E_N-g)=-\int_{[0,1]}r_N(u)\dd\mu_a(u).
\]
The atom at one in \eqref{eq:measure} is killed by $r_N(1)=0$, so
\eqref{eq:criticalexact} follows. Its positive density proves strict
decrease in $N$, and $\nu_a>1$ proves $A_N>I_N$. Since $E_1=0$,
$A_1=-2g$. Corollary~\ref{cor:endbounds} now proves the upper bound.
At $N=1$ and $a\downarrow0$, $A_1(a)\to\CC$, so no smaller
uniform constant is possible.

For fixed $T>0$, the beta law in \eqref{eq:Dv} converges to the point
mass at one as $a\uparrow1$, giving $\nu_a(T)\to1$. Equation
\eqref{eq:nubound} supplies the common integrable majorant
$r_N(u)(1+1/(-\log u))$ for $N\ge1$. Dominated convergence proves
$A_N(a)\to I_N$. The analogous limit as $a\downarrow0$ is $B_N$.
\end{proof}

\begin{proposition}[Exact endpoint recurrences]\label{prop:endrec}
With
\begin{equation}\label{eq:JN}
 J_N=\int_0^1(1-u^2)^N\dd u
 =\frac{4^N(N!)^2}{(2N+1)!},
\end{equation}
one has
\begin{align}
 I_0&=\frac\pi4,& I_{N+1}&=2I_N-J_N\quad(N\ge0),\label{eq:Irec}\\
 B_1&=\CC,& B_{N+1}&=2B_N-J_{N-1}-\frac1{2N}\quad(N\ge1).
 \label{eq:Brec}
\end{align}
Thus every endpoint error is a rational linear combination of $\pi$,
$\log2$, and $1$.
\end{proposition}
\begin{proof}
The beta integral gives \eqref{eq:JN}. The identity
$r_{N+1}=2r_N-(1-u^2)^N$ gives \eqref{eq:Irec}; after division by
$1-u$, it gives \eqref{eq:Brec}. Indeed
$\int_0^1(1+u)(1-u^2)^{N-1}\dd u=J_{N-1}+1/(2N)$.
Finally $B_1=\int_0^1(1+u)/(1+u^2)\dd u=\CC$.
\end{proof}
These recurrences are exact identities, not necessarily stable
floating-point algorithms at large $N$; their homogeneous multiplier
is $2$, while the actual solutions decrease. The rational certificate
uses them only at $N=8$ with ample exact endpoint bounds.

\begin{proposition}[Failure of monotonicity for every truncation]
\label{prop:counterexample}
The function $a\mapsto A_N(a)$ need not decrease. In fact,
\begin{equation}\label{eq:counterexample}
 A_8(1/10)-\lim_{a\downarrow0}A_8(a)>\frac{13}{10000}.
\end{equation}
Thus the true monotonicity theorem for $A_1$ cannot be extended to all
$A_N$.
\end{proposition}
\begin{proof}
The standard-library-only certificate in
\texttt{code/certify\_euler.py} encloses $g_{1/10,9/10}$ using
$160$ finite Euler terms, integer tenth-root bounds, and the proved
critical tail bound. It then combines the exact eight-term Euler
sum with \eqref{eq:Brec}. The resulting rational lower endpoint for
the difference in \eqref{eq:counterexample} is strictly greater than
$13/10000$, checked by integer cross multiplication. For orientation,
\[
 A_8(1/10)=0.359989099661\ldots,\qquad
 B_8=0.358607136933\ldots.
\]
The endpoint limit is established in Theorem~\ref{thm:sharpconstant}.
A decreasing function on $(0,1)$ could not exceed its limit at zero;
this proves the claim. Full rational endpoints, not just these
printed decimals, are in \texttt{data/exact\_certificates.json}.
\end{proof}
This is a counterexample to a \emph{stronger possible extension}, not
a refutation of the predecessor's correctly delimited Conjecture 12.3.

\section{Large-logarithm density identities and uniform remainders}
\label{sec:largeT}
In this section $b=1-a\in(0,1)$ and $f_0(x)=(e^x-1)^{-1}$.
The following formula isolates a cancellation obscured by the two
separate terms in \eqref{eq:oldnew}.

\begin{lemma}[A convergent subtracted Bose integral]\label{lem:subtracted}
For every $T>0$,
\begin{align}\label{eq:subtracted}
 \nu_a(T)-1=\frac1{\BB(a,b)}\bigg\{&
 \int_0^1\frac{r^{b-1}\{(1-r)^{-b}-1\}}{e^{Tr}-1}\dd r\notag\\
 &-\int_1^\infty\frac{r^{b-1}}{e^{Tr}-1}\dd r\bigg\}.
\end{align}
Each of the two displayed integrals is convergent.
\end{lemma}
\begin{proof}
In \eqref{eq:Dv}, put $r=1-v$. The beta average of the excess is
\[
 \frac1{\BB(a,b)}\int_0^1r^{b-1}(1-r)^{-b}
 \left\{f_0(Tr)-\frac1{1-r}f_0\left(\frac{Tr}{1-r}\right)\right\}\dd r.
\]
Keep the difference together and first truncate its lower endpoint
at $\varepsilon>0$. In the second integral substitute
$y=r/(1-r)$; the weight becomes exactly $y^{b-1}\dd y$.
After subtracting the common integral on $(\varepsilon,1)$, the
remaining cutoff mismatch is
$\int_\varepsilon^{\varepsilon/(1-\varepsilon)}
r^{b-1}f_0(Tr)\dd r=O_{b,T}(\varepsilon^b)$, which tends to zero.
The first integrand in \eqref{eq:subtracted} is $O(r^{b-1})$ near
zero, because $(1-r)^{-b}-1=br+O(r^2)$. It is integrable near one
because $b<1$. The second integral has exponential decay.
\end{proof}

\begin{theorem}[Full density expansion]\label{thm:densityasympt}
For fixed $0<a<1$, $b=1-a$, and every integer $M\ge1$,
\begin{equation}\label{eq:densityasympt}
 \nu_a(T)=1+\frac1{\Gamma(a)}
 \sum_{k=1}^M\frac{(b)_k^2}{k!}\zeta(b+k)T^{-b-k}
 +O_{a,M}(T^{-b-M-1})\qquad(T\to\infty).
\end{equation}
Every displayed coefficient is positive. This is an asymptotic
expansion; no claim of convergence of its infinite sum is made.
\end{theorem}
\begin{proof}
Split the first integral in \eqref{eq:subtracted} at $r=1/2$.
On $(0,1/2)$, use
\[
 (1-r)^{-b}-1=\sum_{k=1}^M\frac{(b)_k}{k!}r^k+O_{b,M}(r^{M+1}).
\]
Extend each monomial integral to $(0,\infty)$, making only an
exponentially small error. The second integral and the first
integral over $(1/2,1)$ are exponentially small as well; their
endpoint factors are integrable. For $k\ge1$,
\[
 \int_0^\infty\frac{r^{b+k-1}}{e^{Tr}-1}\dd r
 =\Gamma(b+k)\zeta(b+k)T^{-b-k},
\]
by expansion into positive exponentials and Tonelli. Since
$\BB(a,b)=\Gamma(a)\Gamma(b)$ and
$\Gamma(b+k)=\Gamma(b)(b)_k$, the coefficient is the one stated.
The remainder is bounded by the same Bose integral with $k=M+1$.
\end{proof}

The leading coefficient will be written
\begin{equation}\label{eq:cb}
 c(b)=\frac{b^2\zeta(1+b)}{\Gamma(1-b)}.
\end{equation}
It is positive on $(0,1)$, continuous with zero limits at both
endpoints, and
\begin{equation}\label{eq:cbends}
 c(b)\sim b\quad(b\downarrow0),\qquad
 c(b)\sim\zeta(2)(1-b)\quad(b\uparrow1).
\end{equation}
These follow from the simple pole of $\zeta$ at one and of $\Gamma$
at zero, or directly from their standard integral regularizations.

\begin{theorem}[An explicit parameter-uniform density remainder]
\label{thm:uniformdensity}
For every $0<b<1$ and $T\ge2$,
\begin{equation}\label{eq:uniformdensity}
 \left|\nu_{1-b}(T)-1-c(b)T^{-1-b}\right|
 \le16c(b)T^{-2-b}+4e^{-T/2}.
\end{equation}
The constants do not depend on $b$.
\end{theorem}
\begin{proof}
For $0\le r\le1/2$, Taylor's theorem and $0<b<1$ give
\[
 0\le(1-r)^{-b}-1-br\le8br^2,
\]
because the second derivative is
$b(b+1)(1-r)^{-b-2}\le16b$. After division by $\BB(a,b)$ and
integration, this error is at most
\[
 \frac{8b\Gamma(b+2)\zeta(b+2)}{\Gamma(a)\Gamma(b)}T^{-b-2}
 \le16c(b)T^{-b-2}.
\]
Here $b+1<2$ and $\zeta(b+2)<\zeta(b+1)$.

It remains to bound three tails. The first-integral tail on
$(1/2,1)$ is at most $f_0(T/2)$, since its weight is bounded above
by the full beta density. The omitted leading monomial integral
on $(1/2,\infty)$ is at most
\[
 \frac{b}{\BB(a,b)}\frac{e^{-T/2}}{1-e^{-T/2}}
 \left(\frac{3}{2T}+\frac1{T^2}\right)
 \le f_0(T/2).
\]
We used $r^b\le1+r$, $b/\BB(a,b)<1$, and $T\ge2$.
The second integral in \eqref{eq:subtracted} is bounded by
$e^{-T}/[T(1-e^{-T})]$, since $\BB(a,b)>1$ and $r^{b-1}\le1$
for $r\ge1$. The sum of these three bounds is smaller than
$4e^{-T/2}$ for $T\ge2$. This proves the result.
\end{proof}

\begin{remark}[Cancellation of the apparent first correction]
Equation \eqref{eq:oldnew} contains an explicit $T^{-b}$ term.
The beta average $C_b(T)$ has an opposite contribution of that same
order, so the true excess $\nu_a(T)-1$ begins at $T^{-b-1}$.
Expanding the two old terms separately and retaining only the first
one would give a false Euler-error correction. Equations
\eqref{eq:subtracted} and \eqref{eq:uniformdensity} perform and bound
the cancellation before taking asymptotics.
\end{remark}

\section{Uniform Euler asymptotics and the moving maximum}
\label{sec:uniformEuler}
For $N>1$, write
\begin{equation}\label{eq:ellN}
 \ell_N=\frac12\log N.
\end{equation}
The positive probability measure
$\dd P_N(u)=r_N(u)\dd u/I_N$ localizes at $u$ of order $N^{-1/2}$.
We give the uniform estimates explicitly, because fixed-parameter
asymptotics alone cannot locate a maximizing parameter that moves
with $N$.

\begin{lemma}[The elementary comparison asymptotics]\label{lem:Iasympt}
As $N\to\infty$,
\begin{align}
 I_N&=\frac{\sqrt\pi}{2\sqrt N}
       \left(1-\frac7{8N}+O(N^{-2})\right),\label{eq:Iasympt}\\
 B_N-I_N&=\frac1{2N}+O(N^{-3/2}).\label{eq:Basympt}
\end{align}
\end{lemma}
\begin{proof}
Substitute $u=v/\sqrt N$. On a slowly growing bounded window,
\[
 \frac{(1-v^2/N)^N}{1+v^2/N}
 =e^{-v^2}\left\{1-\frac{v^2+v^4/2}{N}
       +O\left(\frac{v^4+v^6+v^8}{N^2}\right)\right\}.
\]
The complementary tail is controlled by $e^{-v^2}$; enlarging the
window, or splitting at $v=N^{1/12}$, justifies termwise integration
with an $O(N^{-2})$ relative remainder. Gaussian moments give
$\int v^2e^{-v^2}/\int e^{-v^2}=1/2$ and
$\int v^4e^{-v^2}/\int e^{-v^2}=3/4$, yielding $7/8$.
For \eqref{eq:Basympt}, use
$B_N-I_N=\int_0^1u r_N(u)/(1-u)\dd u$. The same substitution gives
the leading integral $N^{-1}\int_0^\infty ve^{-v^2}\dd v=1/(2N)$;
the next term from $(1-u)^{-1}$ is $O(N^{-3/2})$.
\end{proof}

\begin{theorem}[Uniform scaled-error expansion]\label{thm:uniformEuler}
There are absolute constants $C,N_0$ such that, for all $N\ge N_0$
and $0<b<1$,
\begin{equation}\label{eq:uniformEuler}
 \left|\frac{A_N(1-b)}{I_N}-1-c(b)\ell_N^{-1-b}\right|
 \le C\left\{c(b)\ell_N^{-2-b}+N^{-1/4}\right\}.
\end{equation}
The same bound holds at $b=0,1$ by the continuous endpoint extensions
$A_N(1)=I_N$ and $A_N(0)=B_N$, with $c(0)=c(1)=0$.
For each fixed $0<a<1$, in particular,
\begin{equation}\label{eq:fixedEuler}
 E_N-g_{a,1-a}
 =\frac{2^{-N}\sqrt\pi}{2\sqrt N}
 \left\{1+c(1-a)\ell_N^{a-2}
             +O_a(\ell_N^{a-3})\right\}.
\end{equation}
\end{theorem}
\begin{proof}
We average \eqref{eq:uniformdensity} against $P_N$ and control all
exceptional regions uniformly. Put $v=\sqrt N\,u$ and
$T=-\log u=\ell_N-\log v$. On $u\le N^{-1/4}$ one has
$T\ge\ell_N/2$. For $s\in[1,2]$, the mean-value theorem gives
\[
 |T^{-s}-\ell_N^{-s}|
 \le16\ell_N^{-s-1}|\log v|.
\]
By \eqref{eq:Iasympt}, $I_N\ge c_0N^{-1/2}$ for an absolute
$c_0>0$ and all sufficiently large $N$. Also
$r_N(v/\sqrt N)\le e^{-v^2}$. Hence the $P_N$-expectation of
$|\log v|$ is bounded by a constant times
$\int_0^\infty e^{-v^2}|\log v|\dd v<\infty$, uniformly in $N$.
The same estimate shows that the expectation of $v^{1/2}$ is bounded.
It follows, uniformly in $s\in[1,2]$, that
\[
 \E_{P_N}[T^{-s};u\le N^{-1/4}]
 =\ell_N^{-s}+O(\ell_N^{-s-1})
\]
after an exponentially small missing-mass error. The first remainder
in \eqref{eq:uniformdensity} averages to
$O(c(b)\ell_N^{-2-b})$, and the second averages to
$O(\E u^{1/2})=O(N^{-1/4})$.

For the omitted region $u>N^{-1/4}$, use the global inequality
$\nu_a(T)-1<1/T\le1/(1-u)$. The elementary estimate
\[
 \frac{r_N(u)}{1-u}
 =\frac{1+u}{1+u^2}(1-u^2)^{N-1}
 \le2(1-u^2)^{N-1}
\]
bounds its unnormalized contribution by
$2e^{-(N-1)/\sqrt N}$. After normalization it is still smaller
than $O(N^{-1/4})$. This proves \eqref{eq:uniformEuler}.
Endpoint continuity was proved in Theorem~\ref{thm:sharpconstant};
pass to the endpoint limits in the uniform bound.
For fixed $a$, combine \eqref{eq:Iasympt} with
\eqref{eq:uniformEuler}; the inverse-power-in-$N$ errors are smaller
than the displayed logarithmic remainder.
\end{proof}

\begin{theorem}[Asymptotic location and height of the worst parameter]
\label{thm:maximizer}
Extend $A_N$ continuously to $[0,1]$, and let
\begin{equation}\label{eq:UN}
 U_N=\max_{0\le a\le1}A_N(a).
\end{equation}
For every choice of maximizing parameter $a_N$,
\begin{align}
 U_N&=I_N\left\{1+\frac{1+o(1)}{e\,\ell_N\log\ell_N}\right\},
 \label{eq:Uasympt}\\
 (1-a_N)\log\ell_N&\longrightarrow1.\label{eq:amax}
\end{align}
For all sufficiently large $N$, every maximizer lies in $(0,1)$.
The theorem does not assert uniqueness of the maximizer.
\end{theorem}
\begin{proof}
Existence follows from continuity and compactness. Set
$\lambda=\log\ell_N$ and
$M(\lambda)=\max_{0\le b\le1}c(b)e^{-\lambda b}$.
Equation \eqref{eq:cbends} gives
\begin{equation}\label{eq:elementarymax}
 M(\lambda)\sim\frac1{e\lambda}.
\end{equation}
Here is a direct justification. For any fixed $\delta>0$, the
contribution from $b\ge\delta$ is exponentially small in $\lambda$.
For $b<\delta$, $c(b)/b$ approaches one uniformly as $\delta\downarrow0$.
The change of variable $t=\lambda b$ reduces the maximum to that of
$te^{-t}$, whose unique maximum is $1/e$ at $t=1$.

The uniform error in \eqref{eq:uniformEuler} is bounded by
$O(\ell_N^{-2}M(\lambda))+O(N^{-1/4})$, which is
$o(1/(\ell_N\lambda))$. Taking maxima proves \eqref{eq:Uasympt}.
Moreover any maximizing $b_N=1-a_N$ satisfies
$\lambda c(b_N)e^{-\lambda b_N}\to1/e$.
As in the argument above, first $b_N\to0$, and then
$(\lambda b_N)e^{-\lambda b_N}\to1/e$. The uniqueness of the maximum
of $te^{-t}$ gives $\lambda b_N\to1$, proving \eqref{eq:amax}.
Finally $A_N(1)=I_N$, while \eqref{eq:Basympt} implies
$(A_N(0)-I_N)/I_N=O(N^{-1/2})$. Both are smaller than the positive
interior excess in \eqref{eq:Uasympt}, so neither endpoint maximizes
for sufficiently large $N$.
\end{proof}

The shortest truncation and long truncations therefore have different
worst-parameter behavior. At $N=1$ the supremum is approached at
$a=0$. As $N$ grows, the maximizing parameter moves towards $a=1$,
with distance asymptotic to $1/\log(\tfrac12\log N)$. This conclusion
comes from a parameter-uniform estimate, not an extrapolation of the
small numerical grid.

\section{Rational beta profiles and finite polylogarithm identities}
\label{sec:cyclotomic}
The beta identity can be sampled at rational $v$ without introducing
higher-depth polylogarithms. Repeated denominator roots must, however,
be retained.

\begin{theorem}[Finite rational-profile reduction]\label{thm:profile}
Let $0<p<q$ be coprime integers, $w>0$, and $z\in\cut\setminus\{0\}$.
Set
\[
 P=\{\alpha:\alpha^p=z\},\qquad
 Q=\{\alpha:\alpha^q=z\}.
\]
For each $\alpha\in P\cup Q$, define
\begin{equation}\label{eq:profileAB}
 (A_\alpha,B_\alpha)=
 \begin{cases}
 \left(\dfrac{z}{p(\alpha^q-z)},0\right),&\alpha\in P\setminus Q,\\[6pt]
 \left(\dfrac{z}{q(\alpha^p-z)},0\right),&\alpha\in Q\setminus P,\\[6pt]
 \left(-\dfrac{p+q}{2pq},\dfrac1{pq}\right),&\alpha\in P\cap Q.
 \end{cases}
\end{equation}
Then the exact identity is
\begin{equation}\label{eq:profileLi}
 \HH_w(p/q;z)=q^w\sum_{\alpha\in P\cup Q}
 \{A_\alpha\Li_w(\alpha)+B_\alpha\Li_{w-1}(\alpha)\}.
\end{equation}
There is a common root precisely when $z^{q-p}=1$, and in that case
it is unique. All polylogarithm arguments in the formula lie outside
$[1,\infty)$.
\end{theorem}
\begin{proof}
In \eqref{eq:kernel} substitute $t=e^{-T/q}$. This gives
\begin{equation}\label{eq:profileMellin}
 \HH_w(p/q;z)=\frac{q^w}{\Gamma(w)}\int_0^1
 (-\log t)^{w-1}\frac{z^2t^{p+q}}{(1-zt^p)(1-zt^q)}\frac{\dd t}{t}.
\end{equation}
Consider the rational function in this integrand. Its partial
fractions are
\begin{equation}\label{eq:rationalPF}
 \frac{z^2t^{p+q}}{(1-zt^p)(1-zt^q)}
 =\sum_{\alpha\in P\cup Q}
 \left\{A_\alpha\frac{\alpha t}{1-\alpha t}
       +B_\alpha\frac{\alpha t}{(1-\alpha t)^2}\right\}.
\end{equation}
At a simple root, taking the coefficient of $(1-\alpha t)^{-1}$
immediately gives the first two cases of \eqref{eq:profileAB}.
At a common root write $\varepsilon=1-\alpha t$. Then
\[
 \frac{zt^p}{1-zt^p}
 =\frac1{p\varepsilon}-\frac{p+1}{2p}+O(\varepsilon),
\]
and similarly for $q$. Their product has double-pole coefficient
$1/(pq)$ and simple-pole coefficient $-(p+q+2)/(2pq)$.
Since $\alpha t/(1-\alpha t)=\varepsilon^{-1}-1$ and
$\alpha t/(1-\alpha t)^2=\varepsilon^{-2}-\varepsilon^{-1}$,
these are exactly the last pair in \eqref{eq:profileAB}.
The two rational functions in \eqref{eq:rationalPF} now have identical
principal parts at every finite pole. Their difference is a
polynomial bounded at infinity, hence constant; it vanishes at zero.
This proves the identity.

The Mellin integrals of the two basis fractions are $\Li_w(\alpha)$
and $\Li_{w-1}(\alpha)$, respectively. One may prove the second by
differentiating the classical polylogarithm integral with respect to
$\alpha$; this works for $w>0$, including $w\le1$, without using a
divergent boundary series. If $\alpha\ge1$ were real, then
$z=\alpha^p$ or $\alpha^q$ would belong to the excluded cut.
Thus both integrals use the stated principal branches.

Finally a common root satisfies $\alpha^{q-p}=1$, hence
$z^{q-p}=1$. Conversely, if $z^{q-p}=1$, the coprimality
$\gcd(p,q-p)=1$ makes the map $\alpha\mapsto\alpha^p$ a bijection on
the $(q-p)$th roots of unity. This gives exactly one common root.
\end{proof}

\begin{corollary}[Gaussian cyclotomic specialization]\label{cor:cyclo}
At $z=\ii$, all arguments in \eqref{eq:profileLi} are roots of unity
of conductor dividing $4\operatorname{lcm}(p,q)$. A repeated pole
occurs exactly when $4\mid(q-p)$; its root is $\ii$ for $p\equiv1\pmod4$
and $-\ii$ for $p\equiv3\pmod4$.
\end{corollary}
For instance, the repeated-pole case $p=1,q=5$ gives
\begin{align}\label{eq:resonant15}
 \HH_w(1/5;\ii)=5^w\bigg\{&-\frac35\Li_w(\ii)
       +\frac15\Li_{w-1}(\ii)\notag\\
 &+\sum_{\substack{\alpha^5=\ii\\\alpha\ne\ii}}
       \frac{\ii}{5(\alpha-\ii)}\Li_w(\alpha)\bigg\}.
\end{align}
The $\Li_{w-1}(\ii)$ term is forced by a genuine double pole. It cannot
be dropped by applying a simple-root formula at a vanishing denominator.

\begin{proposition}[An elementary half-profile identity]\label{prop:halfprofile}
The weight-one rational profile has the exact Gaussian value
\begin{equation}\label{eq:halfprofile}
 -\Im\HH_1(1/2;\ii)
 =\frac{3\pi}{8}-\frac{\log2}{4}
       -\frac{\log(1+\sqrt2)}{\sqrt2}
 =0.381585209815956\ldots.
\end{equation}
\end{proposition}
\begin{proof}
Equation \eqref{eq:profileMellin} gives the real integral
\[
 2\int_0^1\frac{t^3+t^4}{(1+t^2)(1+t^4)}\dd t.
\]
One antiderivative, verified by differentiation, is
\begin{align*}
 &-\frac12\log(1+t^2)
 +\frac{1+\sqrt2}{4}\log(t^2-\sqrt2t+1)
 +\frac{1-\sqrt2}{4}\log(t^2+\sqrt2t+1)\\
 &\hspace{20mm}+\arctan t
 +\frac12\arctan(\sqrt2t-1)
 -\frac12\arctan(\sqrt2t+1).
\end{align*}
At $t=1$, the two last arctangents are $\pi/8$ and $3\pi/8$;
at $t=0$ they are $-\pi/4$ and $\pi/4$.
Combining the logarithms of $2\pm\sqrt2$ yields \eqref{eq:halfprofile}.
\end{proof}

\begin{remark}[A profile value is not a fractional harmonic value]
The exact identity \eqref{eq:halfprofile} evaluates the deterministic
profile at $v=1/2$. It does \emph{not} evaluate $F_{1/2,1/2}(\ii)$,
which is its nontrivial beta average. Indeed the latter has imaginary
part $-0.4039369877170135798\ldots$, certified below. The distinction
prevents a plausible but false ``midpoint substitution'' identity.
\end{remark}

\subsection{Certified monotone profile quadrature}
Let $J(v)=-\Im\HH_w(v;\ii)$, and choose rational nodes
$0=v_0<\cdots<v_m=1$. For $V\sim\mathrm{Beta}(a,w-a)$ put
$p_j=\Pr(v_{j-1}<V<v_j)$. The strict profile monotonicity proves
\begin{equation}\label{eq:profilebracket}
 \sum_{j=1}^m p_jJ(v_j)
 <-g_{a,w-a}
 <\sum_{j=1}^m p_jJ(v_{j-1}).
\end{equation}
Every interior endpoint profile admits the finite reduction
\eqref{eq:profileLi}; the two extreme profiles are depth-one values
from \eqref{eq:end0}--\eqref{eq:end1}. With certified beta masses and
polylogarithm evaluations, \eqref{eq:profilebracket} therefore gives
another rigorous evaluator. This is a finite quadrature enclosure,
not a finite depth-one reduction of the exact beta average.

\section{Exact certificates, numerical checks, and reproducibility}
\label{sec:verification}

\subsection{Rational power enclosures and the Euler certificates}
The main evaluator is standard-library Python. For a positive integer
$n$, a nonnegative rational exponent $p/q$, and $M=10^D$, it finds an
integer $L$ satisfying
\begin{equation}\label{eq:integerroot}
 L^q n^p\le M^q<(L+1)^q n^p.
\end{equation}
An integer Newton iteration produces the candidate; the two displayed
inequalities are then replayed exactly. Consequently
$n^{-p/q}\in[L/M,(L+1)/M]$, with an exact singleton when equality
holds on the left. Positive interval arithmetic encloses the harmonic
coefficients, and signed interval arithmetic evaluates
\eqref{eq:finiteweights}.

The constants are also enclosed rationally. Machin's identity
$\pi=16\arctan(1/5)-4\arctan(1/239)$, with alternating-series
remainders, bounds $\pi$. The positive expansion
$\log2=2\sum_{k\ge0}[(2k+1)3^{2k+1}]^{-1}$, with a geometric upper
bound on its tail, bounds $\log2$. Thus no floating-point value of
$\CC$ is assumed. Its rational upper endpoint is also checked to be
less than $283/250$.

The default run uses $N=160$ and $D=140$. On the critical line it
subtracts a proved upper bound $\CC 2^{-160}$ from the finite-sum
lower endpoint. For other positive orders it uses
$2(161)2^{-160}$ from Theorem~\ref{thm:allEuler}. The resulting
intervals are recorded as exact fractions. For compact presentation,
Table~\ref{tab:values} gives outward-rounded intervals of width
$10^{-18}$, much wider than the stored rational bounds.

\begin{table}[htbp]\centering\small
\begin{tabular}{@{}ccc@{}}
\toprule
$(a,b)$ & lower bound for $g_{a,b}$ & upper bound for $g_{a,b}$\\
\midrule
$(1/10,9/10)$ & $-0.529761932285771264$ & $-0.529761932285771263$\\
$(1/4,3/4)$   & $-0.479061054846237617$ & $-0.479061054846237616$\\
$(1/2,1/2)$   & $-0.403936987717013580$ & $-0.403936987717013579$\\
$(3/4,1/4)$   & $-0.339754024973484554$ & $-0.339754024973484553$\\
$(9/10,1/10)$ & $-0.306037812999015483$ & $-0.306037812999015482$\\
\midrule
$(1/10,1/10)$ & $-0.490480773154378135$ & $-0.490480773154378134$\\
$(1/4,1/4)$   & $-0.464611276461282331$ & $-0.464611276461282330$\\
$(1/5,1/2)$   & $-0.488913236652073667$ & $-0.488913236652073666$\\
\midrule
$(1/10,1)$ & $-0.530844252966159327$ & $-0.530844252966159326$\\
$(1,1)$ & $-0.272198261287950267$ & $-0.272198261287950266$\\
$(2,1)$ & $-0.113648917599903644$ & $-0.113648917599903643$\\
$(1,2)$ & $-0.251642673723681486$ & $-0.251642673723681485$\\
\bottomrule
\end{tabular}
\caption{Exact rational certificates, displayed with outward decimal rounding.
The middle three rows are below the signed-measure threshold.}
\label{tab:values}
\end{table}

The run checks 1,755 integer-root inequalities and the strict
counterexample \eqref{eq:counterexample}. These are finite certificates
using the analytic tail theorems; they do not replace their proofs.

\subsection{Exact cyclotomic recombination}
The separate standard-library verifier
\texttt{code/verify\_cyclotomic.py} computes in
$\Q[X]/(\Phi_L(X))$, where $L=4\operatorname{lcm}(p,q)$.
It generates $\Phi_L$ by exact polynomial division from $X^L-1$,
uses Euclidean inversion in the number field, and reconstructs
\eqref{eq:rationalPF} after multiplication by the \emph{original}
denominator. Every coefficient is an exact rational vector. It does
not approximate roots of unity or evaluate a polylogarithm.

\begin{center}\small
\begin{tabular}{@{}rrrr@{}}
\toprule
$p/q$ & $L$ & $[\Q(\zeta_L):\Q]$ & repeated pole\\
\midrule
$1/2$&8&4&none\\
$1/3$&12&4&none\\
$1/5$&20&8&$\ii$\\
$2/3$&24&8&none\\
$3/5$&60&16&none\\
$3/7$&84&24&$-\ii$\\
\bottomrule
\end{tabular}
\end{center}
These six cases verify 548 exact rational coordinate equalities.
They are checks of the general partial-fraction proof, not 548
independent relations among periods.

\subsection{Symbolic identities and independent numerical diagnostics}
The SymPy program \texttt{code/verify\_symbolic.py} checks 176 exact
identities: kernel derivatives and decompositions, ten finite Euler
remainders, binomial-tail weights, endpoint recurrences, the Bose
convexity factorization, the half-profile antiderivative, and initial
density Taylor coefficients.

The numerical program \texttt{code/diagnostics.py} is explicitly
separate. At 45 decimal working digits it compares twelve evaluations
of the new density with the independent regularized Hurwitz density,
checks their analytic bounds, and compares four cyclotomic formulas
with direct Mellin quadrature. It also checks the elementary
half-profile and, at 160 working digits, the critical Euler grid and
the classical identity $g_{1,1}=-\pi\log2/8$.
No numerical comparison is used to promote an unproved period identity
to a theorem.

\Needspace{12\baselineskip}
\subsection{Replay and build}
From the report directory, the complete replay commands are
\begin{verbatim}
python code/certify_euler.py
python code/verify_cyclotomic.py
python code/verify_symbolic.py
python code/diagnostics.py
python code/build.py
python code/check_package.py
\end{verbatim}
Only the symbolic and diagnostic programs require optional Python
packages (SymPy and mpmath). The two exact rational verifiers require
only the standard library. The PDF build uses three serial
\texttt{pdflatex} passes. The delivered receipt records source hashes,
build checks, and a rendered review; it is evidence about this artifact,
not a proof-assistant or external-referee endorsement.

\section{Corrections, limits, and repository integration}
\label{sec:audit}

\subsection{Changes justified by the proofs}
The precise editorial change is to promote the predecessor's
\texttt{conj:constant} to a theorem, with a cross-reference to
Theorems~\ref{thm:transfer} and \ref{thm:sharpconstant}. Its endpoint
limits were already established; the new content is the global strict
order comparison and the resulting sharp uniform bound.

The Euler evaluator can now be extended below $a+b=1$, but its proof
must cite Theorem~\ref{thm:allEuler}. It must not cite a finite signed
Hausdorff measure that cannot exist in that region. On and below the
critical line the evaluated Gaussian values remain analytic/Abel
values, not ordinary sums.

For critical large-$N$ analysis, \eqref{eq:subtracted} and
\eqref{eq:uniformdensity} should replace any unproved termwise
large-$T$ heuristic applied to the two separate pieces of the old
density. The cancellation of order $T^{-b}$ is mathematically real.
The counterexample at $N=8$ should accompany any proposed extension
of parameter monotonicity from $A_1$ to all $A_N$.

In profile reductions, the repeated-pole case must be stated
separately; \eqref{eq:resonant15} is a small regression test. A
rational beta-profile value must never be identified with its beta
average by substituting the mean parameter.

\subsection{What is not being corrected or claimed}
No false central theorem in the inspected real-order report was
identified. The new results strengthen it and resolve an explicitly
open conjecture. The warnings above concern invalid extrapolations,
not alleged mistakes that the source did not make. The full collective
manuscript has not been line-by-line re-audited in this continuation.

The $S_4$ proof remains prior material; $S_6$, the broad normalized-radius
conjecture, and below-threshold angular uniqueness are not settled
here. No assertion of minimal depth, linear independence of periods,
or completeness of a relation system follows from the certificates.
There is no Lean formalization in this package.

For bibliographic currency, Liu and Pego's 2016 article has a 2026
corrigendum concerning its Fuss--Catalan Theorem~5 \cite{liupegoerr}.
That result is not used in the present proofs. It is worth recording
in the repository bibliography, but it supplies no reason to retract
the elementary moment arguments proved here.

\subsection{An additive integration plan}
The proposed destination is
\begin{center}\small
\nolinkurl{Analysis/Polylogarithms/docs/reports/beta-transfer-critical-euler/}.
\end{center}
The archive already uses this path. The supplied manuscript fragment
has label prefix \texttt{betatransfer:}. It summarizes the transfer
and Euler results and points to the complete article rather than
silently replacing historical proofs. The status ledger records
exactly which conjecture is resolved, which claims are new theorems,
and which questions remain open.

The preserved predecessor source is identified by Git blob
\nolinkurl{fb8813996c16084014dd3154d4ec299d14996c93} and SHA-256
\nolinkurl{6173220e86b32cd8c23f96216f2e5a2e32813d49fdd7769b62162e15e3cbe63c}.
The retrieved source bytes match that pinned Git blob. No upstream
branch, canonical PDF, or immutable incoming archive was modified in
preparing this delivery.

\section{Further research questions}
\label{sec:future}

\subsection{The exact radius range of order transfer}
The sufficient radius $\sqrt3$ arises from positivity of
$2x+y-\rho^2x^2y$ on the full triangle. At larger radii that pointwise
property fails, but its Mellin integral could still preserve the
order comparison. Determine, for each $w>0$, the maximal radius
range in which \eqref{eq:transfer} holds for every $a\in(0,w)$.
A counterexample must concern the integrated function, not merely a
negative derivative of a single kernel value.

\subsection{Uniqueness and finer drift of the error maximum}
Theorem~\ref{thm:maximizer} locates every maximizing parameter to
leading order but does not prove uniqueness. A concrete conjecture is:
\begin{conjecture}[Eventual single maximum]\label{conj:unimodal}
For all sufficiently large integers $N$, the function $A_N$ has a
unique maximizer in $(0,1)$ and is strictly increasing before it and
strictly decreasing after it.
\end{conjecture}
The present uniform value estimates do not prove this stronger
shape statement. A derivative-uniform version of
\eqref{eq:uniformEuler} would be a natural route. It should also
identify further terms in the drift of $a_N$ and the height of $U_N$.
The first truncation at which an interior maximum appears is a
separate finite problem; the exact $N=8$ counterexample only supplies
an upper bound for the first failure of decreasing behavior.

\subsection{Sharp subcritical Euler asymptotics}
The all-positive-order bound is $O(\sqrt N\,2^{-N})$, whereas the
critical fixed-parameter error is $\asymp2^{-N}N^{-1/2}$.
Determine the sharp asymptotic law for $a+b<1$, including uniform
transition estimates as $a+b\uparrow1$. The divided-difference
kernel gives a valid starting point without postulating a nonexistent
finite signed measure. Its near-diagonal cancellation and its
small-$u$ localization must be analyzed together.

\subsection{Sharper Hurwitz remainders and higher-depth analogues}
The inequalities \eqref{eq:Rderiv} and \eqref{eq:Rdiff} are uniform
but need not have optimal constants at fixed $a$. Can one subtract
additional explicitly positive Hausdorff densities from
\eqref{eq:defectmoments} while retaining a useful all-$q$ theorem?
At higher nested depth, an iterated gamma integral gives more than
two denominators. Determine whether a multivariate order-transfer
principle survives, and whether critical endpoint atoms are replaced
by higher divided differences with controlled signs.

\subsection{Cyclotomic profiles and certified beta quadrature}
Theorem~\ref{thm:profile} reduces each rational profile, not the full
beta average. An efficient certified implementation of
\eqref{eq:profilebracket} should balance rational-node conductor,
beta-mass enclosure, and the steepness of $J(v)$ near the endpoints.
The general-$z$ repeated-pole criterion also suggests a useful
classification of exact identities at different roots of unity.
Any proposed finite reduction of a nontrivial beta average needs an
additional identity, not just finer quadrature.

\subsection{Independent review and formalization}
A compact formalization target is the chain consisting of the
positive gamma integral, the beta covariance lemma, the elementary
Bose convexity inequality, and the finite Euler divided difference.
The exact integer and cyclotomic verifiers can be translated into
formally checked certificates separately. The $S_6$ conjecture and
its existing exact-relation searches remain a distinct research
thread rather than an implied consequence of this analytic work.

\section*{Conclusion}
\addcontentsline{toc}{section}{Conclusion}
A beta distribution exposes a strict order comparison that resolves
the predecessor's critical-constant conjecture. Retaining a positive
divided difference extends Euler evaluation through the signed-measure
threshold. On the boundary, the same framework yields a uniformly
bounded density excess, exact Hurwitz and cyclotomic identities, and
a parameter-uniform asymptotic law whose maximum moves towards the
opposite endpoint from the shortest-truncation extremum. The exact
certificates confirm the stated finite comparisons, while the
remaining global shape and depth questions retain their conjectural
status.

\clearpage
\addcontentsline{toc}{section}{References}
\begin{thebibliography}{9}
\bibitem{repo}
ProveIt Contributors, \emph{Polylogarithms and their Arithmetic Bridges},
canonical manuscript, editorial ledger, and incoming reports,
repository \texttt{VladimirReshetnikov/ProveIt}, revision \pin.
\url{https://github.com/VladimirReshetnikov/ProveIt/tree/5537940fde730713eaac752265a486832b2fe540/Analysis/Polylogarithms/docs/manuscript}.

\bibitem{threshold}
ProveIt Contributors, \emph{The Sharp Real-Order Threshold for Harmonic
Polylogarithms: Endpoint atoms, zero geometry, and certified fractional-order
identities}, 9 October 2026. Preserved article at the same revision,
\nolinkurl{Analysis/Polylogarithms/docs/reports/real-order-threshold/article/real_order_threshold.tex}.
Conjecture 12.3 is the direct target of the present continuation.
\url{https://github.com/VladimirReshetnikov/ProveIt/tree/5537940fde730713eaac752265a486832b2fe540/Analysis/Polylogarithms/docs/reports/real-order-threshold}.

\bibitem{dlmfpolylog}
NIST Digital Library of Mathematical Functions, \S25.12,
\emph{Polylogarithms}, especially the defining series, principal
continuation, and integral representation. Accessed 9 October 2026.
\url{https://dlmf.nist.gov/25.12}.

\bibitem{dlmfhurwitz}
NIST Digital Library of Mathematical Functions, \S25.11,
\emph{Hurwitz Zeta Function}, integral representations and analytic
continuation. Accessed 9 October 2026.
\url{https://dlmf.nist.gov/25.11}.

\bibitem{liupego}
J.-G. Liu and R. L. Pego, \emph{On generating functions of Hausdorff
moment sequences}, Transactions of the American Mathematical Society
\textbf{368} (2016), 8499--8518. DOI: 10.1090/tran/6618.
\url{https://arxiv.org/abs/1401.8052}.

\bibitem{liupegoerr}
J.-G. Liu and R. L. Pego, \emph{Corrigendum to ``On generating
functions of Hausdorff moment sequences''}, Transactions of the
American Mathematical Society \textbf{379} (2026), 3009--3010.
DOI: 10.1090/tran/9582.
The correction concerns the Fuss--Catalan theorem, not a result used
in this article. Author-institution bibliographic record:
\url{https://scholars.duke.edu/publication/1704066}.
\end{thebibliography}
\end{document}
